\uuid{Q5qt}
\titre{Solution générale/condition initiale}
\auteur{hanani ; mzari}
\organisation{exo7}
\licence{CC BY-NC-SA 4.0}
\datecreate{2020-01-20}
\niveau{L1}
\module{Analyse}
\chapitre{Equation différentielle}
\sousChapitre{Résolution d'équation différentielle du deuxième ordre}
\difficulte{3}
\theme{Equations du second ordre, solution générale, condition initiale}
\exo7qcmid{lille-2.equations-differentielles.030}

\contenu{
\begin{question}
\qclassification{225.02}
\qtags{motcle=solution générale/condition initiale}
On considère l'équation différentielle $(E)$ : $y''+y=2\cos x$. Parmi les affirmations suivantes, cocher celles qui sont vraies :
\begin{answers}
\bad{Toute solution de $(E)$ est combinaison linéaire de $\sin x$ et $\cos x$.}
\good{Toute solution de $(E)$ est combinaison linéaire de $\sin x$, $\cos x$, $x\sin x$ et $x\cos x$.}
\good{$(E)$ admet une solution sous la forme $y=x(a\sin x+b\cos x)$, où $a,b,\in \Rr$.}
\bad{La solution de $(E)$ telle que $y(0)=0$ et $y'(0)=0$ est $\displaystyle y=x\cos x$.}
\end{answers}
\begin{explanations}
La solution générale  de l'équation homogène est : $Y=k_1\cos x+k_2\sin x$, $k_1,k_2\in \Rr$. On vérifie que, $\displaystyle y_0=x\sin x $ est une solution particulière de $(E)$. Donc la solution générale de $(E)$ est : $\displaystyle y=k_1\cos x+(k_2+x)\sin x$, pour $k_1,k_2\in \Rr$. Enfin, $y(0)=0\Rightarrow k_1=0$ et $y'(0)=0\Rightarrow k_2=0$.
\end{explanations}
\end{question}
}
